\capitulo{2}{Objetivos del proyecto}

El objetivo general de este proyecto es diseñar e implementar una página web que permita centralizar y optimizar la gestión de horarios, reservas de aulas y planificación académica del campus, mejorando la organización y gestión interna y facilitando la consulta de información tanto a los gestores como al profesorado.


Este proyecto, consiste en diseñar e implementar una aplicación web destinada a las personas demandantes de espacios o aulas y las personas relacionadas con la gestión de reservas en las facultades de la Universidad de Burgos, en concreto se ha limitado en esta primera versión al campus Vena, aunque la arquitectura está preparada para extrapolarlo a otras facultades. En el acceso a la aplicación se distinguirán diferentes roles: administrador, gestor editor, gestor no editor y el personal docente investigador (PDI). Según el rol del usuario que inicie sesión, se le permitirá acceder a unas funcionalidades u otras, según los permisos que disponga.

\section{Objetivos funcionales}

En esta sección se presentan los objetivos funcionales del proyecto, entendidos como aquellas funcionalidades y servicios que la aplicación debe ofrecer para poder satisfacer las necesidades planteadas.

\begin{enumerate}
	\item Gestión de reservas puntuales: Los usuarios autenticados, sin importar su rol, podrán solicitar una reserva indicando la fecha y la hora, el motivo y las características que necesita que tenga el aula. Estas solicitudes serán gestionadas por los gestores, quienes podrán editar los datos de la reserva, y modificar el aula escogido para la reserva. Además, estas reservas dispondrán de un estado (pendiente, aprobada y rechazada) y únicamente el gestor editor podrá crear, modificar o eliminar reservas ya existentes. Cada modificación realizada en ellas, deberá reflejarse automáticamente en la ocupación del aula correspondiente si se viera afectado.
	\item Consulta de aulas y ocupación: El sistema permitirá consultar la ocupación de un aula en un rango de fechas concreto, mostrando tanto las reservas puntuales como las periódicas (ambas o filtrar), y generar una vista semanal basada en las reservas periódicas con el objetivo de crear un informe que poder colocar en la puerta del aula con los detalles de las asignaturas, grupos y horarios.
	\item Consulta de reservas: La aplicación incluye un apartado de consulta de reservas, que permite visualizar y filtrar todas las reservas puntuales registradas en el sistema por responsable, motivo y aula, para facilitar su seguimiento y control.
	\item Gestión de usuarios y roles: Esto será responsabilidad del administrador, que podrá crear, modificar y eliminar usuarios, así como asignarles el rol correspondiente para poder garantizar un control adecuado de los permisos y de las acciones permitidas dentro del sistema. Aunque el proceso de autenticación será delegado a la API institucional del Moodle.
	\item Gestión de horarios y reservas periódicas: Este es uno de los objetivos más importantes del proyecto. El sistema permitirá consultar horarios por grado y curso, así como importar horarios de cursos anteriores como base sobre la que modificar. Permitiendo realizar modificaciones de forma manual sobre el horario, que comprobará si se aplican y cumplen una serie de restricciones que se encontrarán previamente definidas, como la imposibilidad de solapamiento de asignaturas para un mismo grupo, o que se encuentre dentro del horario de apertura de la politécnica. Cuando una modificación vulnere alguna restricción, el sistema mostrará una advertencia, permitiendo al usuario confirmar la acción bajo su responsabilidad y avisando de la violación.
\end{enumerate}

\section{Objetivos técnicos}

A continuación, se presentan los objetivos técnicos del proyecto, entendidos como aquellos objetivos relacionados con la arquitectura, las tecnologías empleadas y las decisiones de implementación necesarias para construir el sistema.

\begin{enumerate}
	\item Diseñar el modelo de bases de datos relacional: se ha diseñado para permitir el almacenamiento de forma estructurada de la información relativa a aulas, grados, asignaturas, reservas, horarios, exámenes y usuarios. Esta estructura permite que se pueda extrapolar a diferentes facultades en un futuro.
	\item Implementar la arquitectura web: su implementación se ha basado en la separación entre el \textit{frontend} y el \textit{backend}, favoreciendo de esta forma la modularidad y mantenibilidad
	\item Implementar sistema de permisos: se ha implementado un sistema de control de acceso basado en roles que limita las acciones disponibles en la aplicación según el tipo de usuario.
	\item Calcular la disponibilidad: que garantiza la actualización de manera automática de la ocupación de las aulas ante cualquier cambio que se realice en las reservas o los horarios.
	\item Incorporar mecanismos de parseo y análisis de documentos existentes: se ha planteado la incorporación de procesos de análisis sobre documentos ya existentes, con el fin de facilitar la carga masiva y la migración de datos al sistema.
	\item Validar restricciones: al realizar una acción se comprueba que no viole ninguna de las restricciones definidas.
	\item Calidad: se realizan validaciones y manejo de errores con el fin de verificar el correcto funcionamiento de la aplicación.
\end{enumerate}



